% TOP_PROBLEM: {"id": 3225, "title": "Strong colorability", "category_id": 3, "category": {"id": 3, "name": "graph_theory", "display_name": "Graph Theory", "description": "Problems involving graphs, networks, and their properties.", "slug": "graph-theory", "order_index": 3, "created_at": "2026-07-31T15:26:25.671Z"}, "status": "open", "problem_number": "OPG-171", "set_id": 12, "statusQualification": "The substantive 2 Delta assertion is stated for Delta at least one; edgeless graphs are handled separately to avoid a zero-colour convention."}
% FIRST_POSTED: 2026-10-01
% ATTEMPT_STATUS: unresolved
% MODEL: GPT 6 Astra Ultra
% UnsolvedMath ID 3225; source catalogue code OPG-171
% !TeX program = lualatex
% Research attempt dated 2026-10-01; canonical statement preserved.
\documentclass[11pt,a4paper]{article}
\usepackage[margin=25mm]{geometry}
\usepackage{fontspec}
\setmainfont{lmroman10-regular.otf}[BoldFont=lmroman10-bold.otf,
 ItalicFont=lmroman10-italic.otf,BoldItalicFont=lmroman10-bolditalic.otf]
\usepackage{amsmath,amssymb,mathtools}
\usepackage{unicode-math}
\setmathfont{latinmodern-math.otf}
\AtBeginDocument{\let\setminus\smallsetminus}
\usepackage{xurl,hyperref}
\hypersetup{colorlinks=true,urlcolor=blue}
\setcounter{secnumdepth}{0}
\setlength{\parindent}{0pt}
\setlength{\parskip}{5pt}
\setlength{\emergencystretch}{3em}
\begin{document}
\section*{Strong colorability}\label{3225}
{\small UnsolvedMath ID 3225 · OPG-171 · Graph Theory · OpenGarden}
\subsection{Definitions and mathematical statement}
Let $G=(V,E)$ be a finite
simple graph and let $r\geq1$
be an integer. A partition
$\mathcal P$ of $V$ into
parts of size at most $r$
consists of nonempty pairwise
disjoint subsets whose union
is $V$. A strong $r$-colouring
relative to $\mathcal P$ is
a function $c:V\to\{1,\ldots,r\}$
that is proper on $G$ and
injective on every part of
$\mathcal P$.

The graph is strongly
$r$-colourable if such a
colouring exists for every
partition into parts of size
at most $r$. Different
partitions may require different
colourings. Equivalently, add
all edges within each part
to make it a clique; every
graph produced this way must
have chromatic number at
most $r$.

The conjecture states that
every finite simple graph
of maximum degree $\Delta\geq1$
is strongly $2\Delta$-colourable.
Explicitly, after any partition
into parts of size at most
$2\Delta$, one must colour
the vertices with $2\Delta$
colours while distinguishing
both adjacent vertices and
vertices sharing a part.

The maximum degree is measured
in $G$ before adding the
partition edges. There is
no restriction on how parts
interact with edges, and no
assumption that $|V|$ is
divisible by $2\Delta$.
The positive-degree condition
avoids the undefined phrase
strongly zero-colourable; edgeless
graphs are trivially strongly
$r$-colourable for every
positive $r$ under this
definition.
\subsection{Short English statement}
After any partition into groups of size at most twice the maximum degree, can a graph be coloured with that many colours so every group also receives distinct colours?
\subsection{Research attempt}
\paragraph{Scope and method.}
This research attempt by GPT-6 Astra Ultra is dated 1 October 2026.
It preserves the catalogue statement and distinguishes elementary proofs
below from cited prior work. No novelty claim or formal proof-assistant
verification is made. The final paragraph states the precise remaining scope.
\paragraph{Literature and a structural approach.}
Axenovich and Martin [R1] prove the conjectured bound for graphs whose
maximum degree is at least one sixth of their order. Their theorem is
prior work; the argument below instead treats arbitrary disjoint unions
of cliques, allowing partitions that cross between those cliques in any
way. We keep the catalogue convention of parts of size at most the
number of colours, rather than require equally sized parts.

\paragraph{Converting two partitions to an incidence multigraph.}
Suppose $G$ is a disjoint union of cliques $C_1,\ldots,C_s$, and let
$P_1,\ldots,P_t$ be the supplied partition. Construct a bipartite
multigraph $H$ with one left vertex for each $C_i$ and one right vertex
for each $P_j$. Each vertex $v$ of $G$ becomes an edge from its clique
to its partition part. Distinct vertices in the same intersection give
parallel edges, which must not be collapsed. A proper edge-colouring of
$H$ gives distinct colours to vertices of $G$ within each $C_i$ and
within each $P_j$. It is therefore exactly a strong colouring of $G$
for this partition. Both implications hold: the incidence construction
does not introduce unwanted constraints between different intersections.

\paragraph{An elementary bipartite edge-colouring proof.}
Let $d$ be the maximum degree of $H$. Add dummy vertices to balance the
two bipartition sizes, then add edges, allowing parallel edges, until
all degrees are $d$. This is possible because total degree deficiency
on the two sides is equal; whenever a deficiency remains choose one
deficient vertex on each side and join them. In the resulting regular
multigraph, every left vertex set $X$ sends $d|X|$ edges to $N(X)$,
whose vertices receive at most $d|N(X)|$ edges. Thus $|N(X)|\ge|X|$.
Hall's theorem yields a perfect matching. Removing it leaves a regular
bipartite graph of degree $d-1$, so induction partitions all edges into
$d$ matchings. Give these matchings different colours and discard the
dummy edges. This proves the required edge-colouring without assuming
anything about simplicity of the incidence graph.

\paragraph{The resulting strong-colouring theorem.}
If $G$ has maximum degree $\Delta$, its cliques have size at most
$\Delta+1$. For every partition with parts of size at most $r$, the
incidence graph has maximum degree at most
$\max(\Delta+1,r)$. Consequently $G$ is strongly $r$-colourable for every
$r\ge\Delta+1$. Taking $r=2\Delta$ proves the catalogue target for this
entire class when $\Delta\ge1$. When $\Delta=1$ every component is a
single vertex or edge, so this also proves the complete maximum-degree-one
case. For an edgeless graph, assign colours independently and injectively
within each part. A clique of order $\Delta+1$ explains why fewer than
$\Delta+1$ colours cannot work universally even before partition edges
are added.

\paragraph{Why two tempting extensions fail.}
For a general graph, adding edges within parts of size $2\Delta$ gives
maximum degree at most $3\Delta-1$. Greedy colouring therefore guarantees
$3\Delta$ colours, not the required $2\Delta$. The incidence proof avoids
that loss by using a genuine clique partition of the original graph;
arbitrary neighbourhoods overlap and do not supply such a partition.

Another strategy chooses an independent transversal, gives it one colour,
and repeats. An arbitrary first transversal can destroy extendibility.
Take one edge $uv$ and isolated vertices $x,y$, with parts $\{u,x\}$ and
$\{v,y\}$. The independent transversal $\{x,y\}$ is legal as a first
colour class. The two remaining vertices are adjacent and cannot both
receive the only remaining colour. A different transversal, such as
$\{u,y\}$, does extend. The failure concerns greedy choice, not existence
of a strong colouring.

\paragraph{Verification and final status.}
The proof explicitly retains repeated intersections as parallel edges and
handles unequal part sizes and isolated vertices. The root batch script
checks the four-vertex failed-transversal example. \textbf{UNRESOLVED.}
Disjoint unions of cliques satisfy a stronger bound, but extending the
matching construction to graphs without a clique-component structure
requires additional control of overlapping adjacency constraints.

\subsection{Sources}
\begin{itemize}
\item[S1] ULAM.AI and the UnsolvedMath Contributors, supplied problems.json, record 3225 (OPG-171), OpenGarden collection. Catalogue identity and source transcription; CC BY 4.0.
\url{https://huggingface.co/datasets/ulamai/UnsolvedMath}.
\item[S2] Open Problem Garden, Strong colorability. Original problem and discussion; the mathematical formulation below is expanded from the supplied source record.
\url{http://www.openproblemgarden.org/op/strong_colorability}.
\item[R1] Maria Axenovich and Ryan R. Martin, \emph{On the strong
chromatic number of graphs}, arXiv:1605.06574. Primary abstract checked
1 October 2026. \url{https://arxiv.org/abs/1605.06574}.
\end{itemize}
\end{document}
